% !TEX root = ../main.tex
\chapter{Удиртгал}\label{ch:intro}

\section{Зорилго}
\emph{Operating System Concepts} номын 4-р бүлгийн “Multithreaded Sorting Application” төслийг Java 17 дээр
хэрэгжүүлэх нь ажлын зорилго юм. Жагсаалтыг хоёр хуваан хоёр sorting thread-ээр эрэмбэлж, гурав дахь merging
thread-ээр хоёр дахь global массив руу нэгтгэнэ. Thread-ийн бодит ажиллагааг хөтөч дээр \emph{Sorting animation}
хэлбэрээр харуулна. Мөн оролтын шалгалт, автомат тест хийж, олон урсгалт шийдэл хэзээ ашигтай болохыг хугацааны
хэмжилтээр шинжилнэ.

\section{Ач холбогдол}
Орчин үеийн компьютер олон цөмтэй боловч нэг урсгалтай програм тэдгээрийн зөвхөн нэгийг ашигладаг. Ажлыг
thread-үүдэд хуваахад хэн ямар өгөгдөлд бичих, хэзээ дуусахыг хүлээх, үр дүнг хэрхэн нэгтгэх гэсэн асуудал
гарч ирдэг~\cite{osc10}. Энэ төсөл нь \emph{өгөгдлийн зэрэгцээ} (data parallelism) болон \emph{fork--join}
загварыг энгийн жишээн дээр харуулдаг.

\section{Хамрах хүрээ}
Програм нь Java 17-ийн зөвхөн стандарт санг ашиглана. Командын мөр (CLI) болон локал web сервер дээрх хөтчийн
UI гэсэн хоёр интерфейстэй. Утгууд \code{int} мужид байх ба хэмжээний хязгаарыг \ref{sec:limits}-р хэсэгт өгөв.
Эх код, тест, гарын авлага GitHub repository-д байгаа~\cite{repo}.
